\documentclass[10pt,a4paper]{amsart}
\usepackage[margin=19mm]{geometry}
\usepackage{amsmath,amssymb,mathtools}
\usepackage[scr=rsfs,cal=euler]{mathalfa}
\usepackage{xfrac}
\usepackage[unicode,hidelinks,bookmarks=false]{hyperref}
\newcommand{\ie}{i.e.\ }
\newcommand{\cf}{cf.\ }
\newcommand{\resp}{respectively\ }
\newcommand{\Subsection}{\subsection}
\providecommand{\nobreakdash}{\nobreak\hbox{-}}
\setlength{\emergencystretch}{2em}
\multlinegap0pt
\usepackage{bm}
\usepackage{bbm}
\usepackage{dsfont}
\usepackage{xspace}
\usepackage{cancel}
\usepackage{mathtools}
\usepackage{amsthm}
\makeatletter
\@ifundefined{theorem}{\newtheorem{theorem}{Theorem}}{}
\@ifundefined{lemma}{\newtheorem{lemma}[theorem]{Lemma}}{}
\@ifundefined{proposition}{\newtheorem{proposition}[theorem]{Proposition}}{}
\makeatother
\title[Concentration Inequalities for Sub-Weibull Random Tensors]{Concentration Inequalities for Sub-Weibull Random Tensors}
\author{Yunfan Zhao}
\date{Source: arXiv:2509.03439v2 (CC BY 4.0)}
\begin{document}
\maketitle

\noindent\emph{Source and licence.} Concentration Inequalities for Sub-Weibull Random Tensors. Version arXiv:2509.03439v2 (CC BY 4.0), distributed under the Creative Commons Attribution 4.0 International licence, as recorded in the arXiv licence metadata for this exact version and shown on the abstract page https://arxiv.org/abs/2509.03439v2. Full rights notice: licenses/arxiv-2509.03439-CC-BY-4.0.txt.\par\smallskip

\noindent\textit{Editorial scope.} Counted unit: the Proposition whose source label is \texttt{prop:4.2} in the article (source0.tex lines 247--256, source label \texttt{prop:4.2}), delivered together with the article's own complete original proof of it (source0.tex lines 258--324, one \texttt{proof} environment of 67 lines). No mathematics is written, repaired or re-derived: statement and proof are the article's own text line for line. Required context retained uncounted: lem:4.1 (lines 224--227). Every definition, display and label of those lines is retained unchanged. Omitted from this package and not redistributed: 1--223, 228--246, 257--257, 325--653. Nothing inside a retained excerpt is omitted or altered: every retained body is delivered exactly as the article's lines are written, so no omission receipt of any kind accompanies this package. Citations: the retained excerpts carry no citation command at all, so no bibliography entry of the article is transcribed and no reference is redistributed. Reference adaptation: none was needed and none was made; every reference inside the delivered unit resolves inside it, so no unresolved reference or citation is produced. Printed slips kept literal: no printed word, symbol, hypothesis or number of the counted statement or of its proof was altered. Layout and class: the article's own class is \texttt{amsart} with packages including amsmath, amssymb, amsthm, mathrsfs, geometry, esint, hyperref, enumitem; the wrapper is the lane's pinned \texttt{amsart} editorial wrapper at 10pt on A4, which restates the theorem-like environments the retained text needs and adds the article's own macro definitions that the retained text needs (none), with extra wrapper packages (bm, bbm, dsfont, xspace, cancel, mathtools). The delivered numbering is the unit's own. Assets: no retained span contains a figure environment, a graphics-inclusion command, a PDF-page inclusion, a table or a TikZ picture; no image file is copied into this package, the package declares no asset dependency and no image file is redistributed with it. Licence: version arXiv:2509.03439v2 of this article is distributed under the Creative Commons Attribution 4.0 International licence, as recorded in the arXiv licence metadata for this exact version and shown on the abstract page https://arxiv.org/abs/2509.03439v2. A full rights notice is delivered at licenses/arxiv-2509.03439-CC-BY-4.0.txt. Characters: every retained excerpt is pure ASCII, so the delivered engine, XeLaTeX with its default Latin Modern text font, reports no missing character.
\begin{lemma}\label{lem:4.1}
Let $n \in \mathbb{N}$ and $\alpha \in [1, 2]$. Let $x$ be a random vector in $\mathbb{R}^n$ whose coordinates $(x_i)_{1 \leq i \leq n}$ are independent, centered (meaning $\mathbb{E}[x_i]=0$), and unit-variance (meaning $\mathbb{E}[x_i^2]=1$) random variables belonging to the class $\mathcal{S}_\alpha$ with sub-exponential norm bounded by a constant $K$. Then, there exists a constant $c > 0$ depending only on $\alpha$ and $K$ such that for all $t \geq 0$:
\begin{equation} \mathbb{P}\left( \left| \|x\|_2 - \sqrt{n} \right| \geq t \right) \leq 2 \exp\left( -c \min\left( t^2, t^\alpha \right) \right). \end{equation}
\end{lemma}

\begin{proposition}[The Generalized Maximal Inequality]\label{prop:4.2}
Let $X = x_1 \otimes \dots \otimes x_d$ be a simple random tensor of class $\mathcal{S}_{\alpha}$. Define the event $\mathcal{E}$ (the "Good Event") as:
\begin{equation}
\mathcal{E} := \bigcap_{k=1}^d \left\{ \prod_{j \neq k} \|x_j\|_2 \le C_0 n^{(d-1)/2} \right\}
\end{equation}
where $C_0 > 1$ is a fixed constant. There exists a threshold $d_0(n, \alpha)$ such that for all $d \le d_0$:
\begin{equation}
\mathbb{P}(\mathcal{E}^c) \le 2d \exp(-c n^{\alpha/2}).
\end{equation}
\end{proposition}

\begin{proof}
We begin by analyzing the probability of the complement event $\mathcal{E}^c$. By De Morgan's laws, the complement is the union of the failure events for each mode $k$:
\begin{equation}
\mathcal{E}^c = \bigcup_{k=1}^d \mathcal{A}_k, \quad \text{where } \mathcal{A}_k := \left\{ \prod_{j \neq k} \|x_j\|_2 > C_0 n^{(d-1)/2} \right\}.
\end{equation}
Applying the subadditivity of the probability measure (the union bound), we obtain:
\begin{equation}
\mathbb{P}(\mathcal{E}^c) \le \sum_{k=1}^d \mathbb{P}(\mathcal{A}_k).
\end{equation}
Since the random vectors $x_1, \dots, x_d$ are independent and identically distributed (in terms of their norm properties within the class $\mathcal{S}_{\alpha}$), the probabilities $\mathbb{P}(\mathcal{A}_k)$ are identical for all $k$. Thus, it suffices to bound the probability of the event $\mathcal{A}_1$, as:
\begin{equation}
\mathbb{P}(\mathcal{E}^c) \le d \cdot \mathbb{P}(\mathcal{A}_1).
\end{equation}

Consider the event $\mathcal{A}_1$:
\begin{equation}
\mathcal{A}_1 = \left\{ \prod_{j=2}^d \|x_j\|_2 > C_0 n^{(d-1)/2} \right\}.
\end{equation}
We normalize the inequality by dividing by $n^{(d-1)/2}$:
\begin{equation}
\mathcal{A}_1 = \left\{ \prod_{j=2}^d \frac{\|x_j\|_2}{\sqrt{n}} > C_0 \right\}.
\end{equation}
Applying the Inequality of Arithmetic and Geometric Means (AM-GM) to the scalars $Y_j := \frac{\|x_j\|_2}{\sqrt{n}} \ge 0$, we have:
\begin{equation}
\left( \prod_{j=2}^d Y_j \right)^{\frac{1}{d-1}} \le \frac{1}{d-1} \sum_{j=2}^d Y_j.
\end{equation}
Consequently, the event $\mathcal{A}_1$ implies the following condition on the sum:
\begin{equation}
\frac{1}{d-1} \sum_{j=2}^d Y_j > C_0^{\frac{1}{d-1}}.
\end{equation}
Let $\delta := C_0^{\frac{1}{d-1}} - 1$. Since $C_0 > 1$, we have $\delta > 0$. The implication becomes:
\begin{equation}
\mathcal{A}_1 \subseteq \left\{ \sum_{j=2}^d \left( \frac{\|x_j\|_2}{\sqrt{n}} - 1 \right) > (d-1)\delta \right\} = \left\{ \sum_{j=2}^d (\|x_j\|_2 - \sqrt{n}) > (d-1)\delta \sqrt{n} \right\}.
\end{equation}

We recall Lemma~\ref{lem:4.1} ($\alpha$-Concentration of the Euclidean Norm), which states that for $\tau > 0$:
\begin{equation}
\mathbb{P}(|\|x\|_2^2 - n| > \tau) \le 2 \exp\left( -c \frac{\tau^\alpha}{n^{\alpha/2}} \right).
\end{equation}
We derive a concentration bound for the linear deviation $\|x\|_2 - \sqrt{n}$. Observe the algebraic identity $|\|x\|_2^2 - n| = |\|x\|_2 - \sqrt{n}| \cdot (\|x\|_2 + \sqrt{n})$. For a deviation of the form $\|x\|_2 - \sqrt{n} > u$ (where $u > 0$), we have $\|x\|_2 + \sqrt{n} > 2\sqrt{n}$. Thus implies $\|x\|_2^2 - n > 2u\sqrt{n}$. Substituting $\tau = 2u\sqrt{n}$ into Lemma~\ref{lem:4.1} yields:
\begin{equation}
\mathbb{P}(\|x\|_2 - \sqrt{n} > u) \le \mathbb{P}(\|x\|_2^2 - n > 2u\sqrt{n}) \le 2 \exp\left( -c \frac{(2u\sqrt{n})^\alpha}{n^{\alpha/2}} \right).
\end{equation}
Simplifying the exponent gives $\frac{(2u\sqrt{n})^\alpha}{n^{\alpha/2}} = 2^\alpha u^\alpha$. Thus, the variable $\xi_j := \|x_j\|_2 - \sqrt{n}$ exhibits $\alpha$-subexponential tail decay with a rate independent of $n$:
\begin{equation}
\mathbb{P}(\xi_j > u) \le 2 \exp(-c' u^\alpha).
\end{equation}

We must bound the probability of the event $\mathcal{B} := \{ \sum_{j=2}^d \xi_j > K \sqrt{n} \}$, where $K = (d-1)\delta$. For independent random variables with $\alpha$-heavy tails ($\alpha \in (0, 2]$), the tail probability of the sum for large deviations is dominated by the probability that a single variable takes a large value. Utilizing the union bound on the event that any single $\xi_j$ exceeds the threshold, and observing that for $d \ll \sqrt{n}$ (ensured by the threshold $d \le d_0$), the term $K\sqrt{n}$ is in the large deviation regime. Substituting $u \approx \sqrt{n}$ into the single variable bound derived above:
\begin{equation}
\mathbb{P}(\xi_j > c \sqrt{n}) \le 2 \exp(-c' (\sqrt{n})^\alpha) = 2 \exp(-c' n^{\alpha/2}).
\end{equation}
Summing over $j=2, \dots, d$:
\begin{equation}
\mathbb{P}(\mathcal{A}_1) \le (d-1) \cdot 2 \exp(-c n^{\alpha/2}) \le 2d \exp(-c n^{\alpha/2}).
\end{equation}

Returning to the initial union bound, we have $\mathbb{P}(\mathcal{E}^c) \le d \cdot \mathbb{P}(\mathcal{A}_1)$. Substituting the bound for $\mathbb{P}(\mathcal{A}_1)$ yields:
\begin{equation}
\mathbb{P}(\mathcal{E}^c) \le 2d^2 \exp(-c n^{\alpha/2}).
\end{equation}
Absorbing the polynomial factor in $d$ into the exponential constant (valid for sufficiently large $n$), we obtain the stated bound:
\begin{equation}
\mathbb{P}(\mathcal{E}^c) \le 2d \exp(-c n^{\alpha/2}).
\end{equation}
This completes the proof.
\end{proof}

\end{document}
